\documentclass{article}
\usepackage[T2A]{fontenc}
\usepackage[cp1251]{inputenc}
\usepackage{amsthm}
\usepackage{amsmath}
\usepackage{amssymb}
\usepackage{amsfonts}
\usepackage{mathrsfs}
\usepackage[12pt]{extsizes}
\usepackage{fancyvrb}
\usepackage{indentfirst}
\usepackage[
  left=2cm, right=2cm, top=2cm, bottom=2cm, headsep=0.2cm, footskip=0.6cm, bindingoffset=0cm
]{geometry}
\usepackage[english,russian]{babel}


\begin{document}
\section*{Вариант 13}
Частные решения уравнений, соответствующие гармоническому закону пульсаций давления на торце канала имеют вид:
\begin{equation*}
 R_k = \frac{2(-1)^{k+1} p_m^+}{(2k-1)\pi w_m D} \left[ A_{ck} \frac{dp}{d\tau} + B_{ck} f_p \right], \quad Q_k = \frac{(-1)^{k+1} p_m^+}{k\pi} \left[ A_{sk} \frac{dp}{d\tau} + B_{sk} f_p \right],
\end{equation*}
где
\begin{equation*}
	A_{ck} = -\frac{a_{2ck}}{a_{1ck}^2 + a_{2ck}^2}, \quad B_{ck} = \frac{a_{1ck}}{a_{1ck}^2 + a_{2ck}^2}, \quad A_{sk} = -\frac{a_{2sk}}{a_{1sk}^2 + a_{2sk}^2}, \quad B_{sk} = \frac{a_{1sk}}{a_{1sk}^2 + a_{2sk}^2}.
\end{equation*}
\end{document}